\documentclass{article}

\title{Problem Incident Review INC0463342}
\author{Stuart Cooper\\stuart.cooper@uts.edu.au\\University of
  Technology, Sydney}
\date{October 9, 2026}

% keep the LaTeX as dead simple as possible as I'd like to do a
% perlpod version of this and allow reformat to HTML, txt, man, LaTeX
% you name it. But to be honest the first version of this is in Ye
% Olde LaTeX

\begin{document}

\maketitle

\section{Root Causes}

\subsection{Assumptions}

%FIXME: itemize three or four assumptions??

Simplistic assumptions made the work seem a lot simpler than it turned
out to be. It was assumed that the changes at the Linux side would be
entirely at the DNS level; change the MX (Mail eXchanger) records for
several domains like EXAMPLE1.uts.edu.au and EXAMPLE2.uts.edu.au,
watch those changes propagate to the Internet and go home on a
Saturday afternoon.

It was also assumed that Abnormal was a drop-in replacement for
Mimecast. Whatever Mimecast had done for us in the past, Abnormal
would do after a cutover.

\subsection{Changes too early}

The DNS changes just mentioned were done on August 22nd while Mimecast
service was switched off on September FIXME. It's normally a great
idea to get part of the work done well ahead of time, but in this
instance a five-week gap between the DNS changes and the Mimecast
switchoff meant a lot of effort was needed to get back into an
email-fix frame of mind and the August changes were not in our recent
memory.

\subsection{No Mimecast involvement}

I'm guessing that any support contract UTS had with Mimecast had
expired well before service was switched off and it would have been
expensive and even embarassing to engage their support just as we
were removing them from the UTS email ecosystem. The fact remains that
we were unsure just what Mimecast did and had to discover the
behaviour in an incident situation.

\section{What went wrong}

I was genuinely surprised that the initial focus of the incidents was
on the listserv issue and not outgoing email to UTS and the Internet
in general. Why were tickets on failed email lists growing in number
while nobody was reporting the (to my mind) larger issue of missing
emails?

It turns out the problem of missing outgoing emails was reported, but
ServiceNow automatically closed the ticket. A Service Now analyst
correctly asked the reported for more information, this information
was provided, but the ticket stayed in ``Customer On-Hold'' status for
three days and was auto-closed. The issue was not surfaced to the
Linux team until FIXME.

%FIXME: explanation of weekly Tuesday supplier invoice payments on
%Tuesdays

\section{What went right}

The MIM Teams channel performed well and there was excellent
co-operation across a range of disciplines featuring Microsoft Email
(Taylor Johnson), Networking and F5 (David Garabito), Site
Certificates (Felix Wu) and Identity (Iain Rayner).

The MIM bridge on the listserv issue was led by FIXME and the
following week's MIM bridge on outgoing email was led by FIXME.
Both leaders did a fine job and gave the technical workers time and
capacity to develop fixes.

The postfix email program is a best-of-breed solution for Linux (and
before then, UNIX) systems. While running older versions on old
versions of RedHat Enterprise Linux, it handled client-side and
server-side certificate management with a simple configuration change.
Postfix is easy to configure and during the incident we were able to
apply changes and have them in effect in the space of minutes.

The GNU mailman program is the earliest open source program still in
widespread use to be written in Python. 

\section{Esential Recommendations}

I offer three recommendations that will help prevent issues the next
time there is a major email change.

\subsection{Document email flows on a single Confluence page}

All end-to-end email flows need to be documented on a single
Confluence page.

Internal Linux to UTS emails
Internal Linux to Internet emails
Internal Linux to listserv emails
% FIXME: is Internal Linux to neomail a thing??
Internal Linux to neomail

\subsection{Documentation of Abnormal}

A separate Confluence page should explain the role of Abnormal in the
UTS email ecosystem.

\subsection{Establish standard email addresses}

The standard email addresses hostmaster@uts.edu.au (for Internet users
to report DNS issues) aud postmaster@uts.edu.au (for Internet users to
report email issues) should be set up. These email addresses can be
done as Microsoft DLs (Distribution Lists), exclusively to @uts.edu.au
adddresses, that grow and shrink over time as staff join and leave.

\section{Further recommendations}

These recommendations should occur through normal quarterly planning
processes.

\subsection{Replace listserv email lists with Microsoft DLs}

While GNU mailman does its job well, more people in an enterprise will
be familiar with Microsoft Distribution Lists. Indeed, some listserv
mailing lists were already migrated to DLs by Suggi and Jorge. Mailman
provides a web frontend to messages which DLs probably don't.

\subsection{Upgrade Linux postoffice servers and postfix versions}

The two postoffice Linux servers should be upgraded to RHEL 9 and the
standard postfix version provided with that version of RedHat
Enterprise Linux.

\subsection{Review the role of cluebringer}

A large Perl application, cluebringer, is used in the UTS mail
ecosystem for rate limiting emails and other functions. 


\end{document}
